% !TeX program = lualatex
\documentclass[11pt, a4paper]{report}

% --- UNIVERSAL PREAMBLE BLOCK ---
\usepackage[a4paper, top=2.5cm, bottom=2.5cm, left=2.5cm, right=2.5cm]{geometry}
\usepackage{fontspec}

\usepackage[spanish, bidi=basic, provide=*]{babel}

\babelprovide[import, onchar=ids fonts]{spanish}
\babelprovide[import, onchar=ids fonts]{english}

% Set default/Latin font to Serif in the main (rm) slot for an academic look
\babelfont{rm}{Noto Serif}
% Assign a specific font for Spanish text
\babelfont[spanish]{rm}{Noto Serif}

% Add because main language is not English
\usepackage{enumitem}
\setlist[itemize]{label=-}
% --- END UNIVERSAL PREAMBLE BLOCK ---

\usepackage{amsmath}
\usepackage{booktabs}
\usepackage{titlesec}
\usepackage{setspace}

% Formato de los capítulos para que sean más estéticos
\titleformat{\chapter}[display]
  {\normalfont\bfseries\Huge}{\chaptertitlename\ \thechapter}{20pt}{\Huge}

% Hyperref siempre al final
\usepackage[colorlinks=true, linkcolor=blue, citecolor=black, urlcolor=blue]{hyperref}

\begin{document}

\begin{titlepage}
    \centering
    \vspace*{2cm}
    
    {\Large \textbf{TRABAJO DE FIN DE GRADO}} \\[1.5cm]
    
    {\huge \textbf{Orchestration Trade-offs in Edge Computing: Systemd vs Kubernetes}}\\[0.5cm]
    {\Large Análisis de rendimiento, distribución de cargas e inferencia de modelos de IA en el extremo de la red}\\[2cm]
    
    {\Large \textbf{Autor:}} \\
    {\large [Tu Nombre y Apellidos]}\\[1cm]
    
    {\Large \textbf{Director/es:}} \\
    {\large [Nombre del Director/a]}\\[2cm]
    
    {\large \textbf{Grado en Ingeniería [Informática / Telecomunicaciones]}} \\[0.5cm]
    {\large [Nombre de tu Universidad]} \\[0.5cm]
    {\large [Mes y Año de la convocatoria]}
    
    \vfill
\end{titlepage}

\pagenumbering{roman}
\tableofcontents
\listoffigures
\listoftables
\cleardoublepage

\pagenumbering{arabic}
\setstretch{1.15}

\chapter{Introducción}

\section{Contexto y Motivación}
El auge del \textit{Edge Computing} plantea nuevos retos en la orquestación de servicios descentralizados. Los entornos perimetrales poseen recursos limitados que cuestionan la viabilidad de modelos \textit{cloud-native} puros.

\section{Objetivos del Trabajo}
El objetivo principal es analizar y comparar los \textit{trade-offs} (rendimiento, consumo y complejidad) entre una orquestación basada en Systemd (nativa y ligera) frente a Kubernetes (robusta pero pesada) en escenarios Edge.

\section{Estructura de la Memoria}
El documento se estructura siguiendo la progresión del desarrollo técnico iterativo, desde la base teórica hasta la comparativa final.

\chapter{Estado del Arte y Tecnologías Involucradas}

Este capítulo compila la investigación técnica previa (correspondiente a las \textit{Setmanas 4 y 5}) necesaria para fundamentar el diseño arquitectónico del proyecto.

\section{Paradigmas de Computación: Cloud vs. Edge Computing}
Definición de conceptos y necesidades de orquestación en el extremo.

\section{Tecnologías de Contenedorización Ligera}
Análisis de Podman frente a Docker, destacando su arquitectura \textit{daemonless} y ejecución \textit{rootless}, ideales para entornos perimetrales seguros y con pocos recursos.

\section{Orquestación y Ciclo de Vida en el Edge}
\subsection{Systemd y la evolución hacia Quadlets}
Gestión del ciclo de vida de los contenedores de forma nativa desde el sistema operativo.
\subsection{Distribuciones Ligeras de Kubernetes (K3s)}
Concepto de K3s como estándar de facto para llevar la semántica de Kubernetes al Edge.

\section{Protocolos de Comunicación y Serialización para Alta Frecuencia}
\subsection{Arquitecturas HTTP/REST y RPC (gRPC)}
\subsection{Mensajería Asíncrona y Punto a Punto (ZeroMQ, MQTT)}
\subsection{Estrategias de Serialización: JSON vs Protobuf vs MessagePack}
Comparativa teórica entre texto plano con alto \textit{overhead} y esquemas binarios eficientes.

\section{Frameworks de Machine Learning para Inferencia}
Introducción a PyTorch y las características del dataset MNIST utilizado para las cargas de trabajo.

\chapter{Metodología y Evolución de la Arquitectura}

\section{Enfoque de Desarrollo Iterativo}
El proyecto se ha abordado mediante iteraciones incrementales semanales. Se partió del aprovisionamiento base (\textit{Setmana 0 a 3}), pasando por la optimización de protocolos (\textit{Setmana 4}), la integración de IA (\textit{Setmana 5}) y culminando en la migración a Kubernetes (\textit{Setmana 6}).

\section{Definición del Entorno de Pruebas}
Infraestructura compuesta por un nodo Master (Ubuntu 22.04, WSL) y nodos Workers (Ubuntu 24.04, VMware VMs).

\section{Herramientas de Aprovisionamiento y Automatización}
Justificación del uso de Ansible para garantizar despliegues reproducibles bajo el patrón de \textit{Side-loading} (evitando repositorios de contenedores externos).

\section{Casos de Uso y Definición de Tareas}
\subsection{Tarea 1: Transmisión masiva de datos (Conteo de imágenes)}
División del dataset en $N$ particiones y retorno del conteo.
\subsection{Tarea 2: Inferencia distribuida de Inteligencia Artificial}
Sustitución del conteo simple por la predicción numérica usando modelos pre-entrenados.

\chapter{Fase 1: Arquitectura Base y Optimización de Red (Systemd + Podman)}

\section{Despliegue de la Infraestructura mediante Systemd Quadlets}
Explicación de cómo se estructuraron los servicios base de forma agnóstica a la carga de trabajo.

\section{Diseño del Banco de Pruebas de Comunicaciones}
Metodología para medir latencias y consumos enviando 60.000 imágenes a los nodos.

\section{Resultados: Comparativa de Rendimiento de Protocolos}
Análisis del rendimiento base (HTTP/1.1), el estándar microservicios (gRPC), el crudo P2P (ZeroMQ) y el paradigma IoT (MQTT).

\section{Resultados: Impacto de la Serialización}
La transición de JSON a binario (Protobuf) y el refinamiento final con MessagePack (\textit{schema-less}).

\section{Selección de la Arquitectura Óptima para el Edge}
Conclusión de esta fase: La elección de ZeroMQ + MessagePack como pila de red ganadora para integrar en Systemd.

\chapter{Fase 2: Implementación del Flujo de Machine Learning}

\section{Entrenamiento y Empaquetado del Modelo (PyTorch)}
El nodo Master asume la carga computacional de entrenamiento y serializa el modelo (`.pth`).

\section{Distribución Segura de Modelos y Datos a Nodos Trabajadores}
Envío del modelo y las particiones mediante la arquitectura de red optimizada en la Fase 1.

\section{Ejecución de Inferencia Descentralizada}
Los Workers cargan el modelo en memoria y evalúan cada imagen, devolviendo las predicciones agregadas al Master.

\section{Análisis de Rendimiento y Consumo Computacional}
Evaluación del incremento de estrés en CPU y memoria ($T_{proc}$) al añadir el \textit{overhead} matemático de los tensores.

\chapter{Fase 3: Migración a Kubernetes y Análisis Comparativo Final}

\section{Transición de Arquitectura: De Quadlets a Manifiestos de K3s}
Adaptación de los servicios a la semántica de \textit{Deployments} y \textit{Services} de K8s.

\section{Despliegue del Escenario de Inferencia sobre K3s}
Ejecución de la Tarea 2 a través de la interfaz de red de contenedores (CNI) de Kubernetes.

\section{Comparativa de Trade-offs: Systemd vs. Kubernetes}
\subsection{Plano de Gestión: Consumo de RAM y CPU}
Comparativa del \textit{overhead} en reposo del Kubelet frente al demonio de Systemd.
\subsection{Plano de Datos: Latencia de red y Throughput}
Impacto del enrutamiento de Kube-proxy frente a la red nativa punto a punto.
\subsection{Complejidad Operativa y Facilidad de Mantenimiento}

\chapter{Conclusiones y Trabajo Futuro}

\section{Conclusiones de la Investigación}
Resumen de los hitos alcanzados y la viabilidad técnica demostrada a lo largo de los capítulos.

\section{Respuesta al Dilema: ¿Systemd o Kubernetes para el Edge?}
Veredicto final apoyado en las métricas extraídas en la Fase 3. Cuándo merece la pena asumir el coste de Kubernetes y cuándo es preferible la simplicidad de Systemd.

\section{Líneas de Trabajo Futuro}

\vspace{1cm}

\begin{thebibliography}{99}
\bibitem{podman_docs} Red Hat, \textit{Podman Documentation}, [Online]. Disponible en: \url{https://podman.io/docs/}
\bibitem{systemd_quadlets} Freedesktop.org, \textit{systemd.service — Service unit configuration}, [Online].
\bibitem{grpc_protocol} Cloud Native Computing Foundation (CNCF), \textit{gRPC: A high performance, open source universal RPC framework}, [Online].
\bibitem{zeromq_guide} P. Hintjens, \textit{ZeroMQ: Messaging for Many Applications}, O'Reilly Media, 2013.
\bibitem{k3s_arch} Rancher Labs (SUSE), \textit{K3s Lightweight Kubernetes}, [Online]. Disponible en: \url{https://k3s.io/}
\bibitem{pytorch_edge} Paszke, A. et al., \textit{PyTorch: An Imperative Style, High-Performance Deep Learning Library}, NeurIPS 2019.
\end{thebibliography}

\end{document}